\usepackage{xeCJK}
\usepackage{geometry}
\geometry{left=2cm,right=2cm,top=2.5cm,bottom=2.5cm}
\setlength{\headheight}{12.7pt}

\usepackage{titlesec}
\titleformat{\section}
  {\normalfont\large\bfseries}
  {\thesection}
  {1em}
  {}
\titlespacing*{\section}
  {0pt}
  {3.5ex plus 1ex minus .2ex}
  {2.3ex plus .2ex}

\usepackage{amsmath}
\usepackage{amssymb}
\usepackage{amsthm}
\usepackage{enumitem}
\setlist[enumerate]{itemsep=0pt,topsep=0pt,parsep=0pt,leftmargin=0pt,itemindent=2.5em}
\setlist[itemize]{itemsep=0pt,topsep=0pt,parsep=0pt,leftmargin=0pt,itemindent=2.5em}
\usepackage{fancyhdr}
\pagestyle{fancy}
\fancyhf{}
\usepackage{graphicx}
\usepackage{hyperref}
\usepackage{indentfirst}
\usepackage{lmodern}


\begin{document}
\setlength{\abovedisplayskip}{1pt}
\setlength{\belowdisplayskip}{1pt}

\section{基数乘法}

\hypertarget{sec:定义1.1}{}
\noindent\textbf{定义1.1 (基数乘法)}

对任意的\href{../01 基数/01 基数#nameddest=sec:定义1.1}{基数}
$\kappa,\lambda$, 定义
\[
    \kappa\otimes\lambda\overset{\mathrm{def}}{=}|\,\kappa\times\lambda\,|,
\]
称之为\textbf{基数乘法}.

\noindent{\kaishu 备注.
这里为了与序数乘法区分, 使用$\otimes$符号. 在不产生歧义的情况下仍然使用$\cdot$符号.}

\vspace{10pt}

显然$\kappa\otimes\lambda=|\,\kappa\lambda\,|$.

\vspace{10pt}

在基数乘法的规律之前, 先给出重要的Hessenberg定理.

\vspace{10pt}

\hypertarget{sec:定理1.1}{}
\noindent\textbf{定理1.1 (Hessenberg定理)}

任给无限序数$\alpha$, 有$\alpha\times\alpha\approx\alpha$.

\noindent{\kaishu 证明.}

对无限序数$\alpha$归纳证明$\alpha\times\alpha\approx\alpha$.
任取无限序数$\alpha$, 假设对任意的无限序数$\beta<\alpha$
有$\beta\times\beta\approx\beta$.

如果$\alpha$不是基数, 则$\kappa=|\alpha|<\alpha$, 从而依假设有
\[
    \alpha\times\alpha\approx\kappa\times\kappa\approx\kappa\approx\alpha.
\]

如果$\alpha$是基数, 那么定义映射
\[
    f:\alpha\times\alpha\to\alpha\times\alpha\times\alpha,\quad
    (x,y)\mapsto(\,\max\{x,y\},\,x,\,y\,).
\]
显然$f$是单射, 把值域上的字典序复制到$\alpha\times\alpha$上,
也记为$<_{\mathrm{lex}}$.

对任意的$\beta<\alpha$, 任取$(x,y)\in\beta\times\beta$和
$(x',y')<_{\mathrm{lex}}(x,y)$, 有
\vspace{3pt}
\[
    \max\{x',y'\}\leqslant\max\{x,y\}<\beta
    \quad\implies\quad(x',y')\in\beta\times\beta,\\[3pt]
\]
即$\beta\times\beta$是前段. 此时:
\begin{enumerate}
    \item 如果$\beta$有限,
    则$\beta\times\beta\approx\beta^2\prec\omega\preccurlyeq\alpha$,
    \item 如果$\beta$无限, 则依假设有
    $\beta\times\beta\approx\beta\prec\alpha$,
\end{enumerate}
综上有$\mathrm{Type}(\beta\times\beta)\prec\alpha$,
即$\mathrm{Type}(\beta\times\beta)<\alpha$.
另外显然$\displaystyle\alpha\times\alpha=
\bigcup_{\beta<\alpha}\beta\times\beta$, 所以
\vspace{3pt}
\[
    \mathrm{Type}(\alpha\times\alpha)=
    \sup_{\beta<\alpha}\mathrm{Type}(\beta\times\beta)\leqslant\alpha.
\]
从而$\alpha\times\alpha\preccurlyeq\alpha$,
显然$\alpha\times\alpha\approx\alpha$. \hfill $\qedsymbol$





\newpage

\hypertarget{sec:定理1.2}{}
\noindent\textbf{定理1.2}

任给非零基数$\kappa,\lambda$,
\begin{enumerate}
    \item 如果$\kappa,\lambda$都有限, 那么$\kappa\otimes\lambda=\kappa\lambda$,
    \item 如果$\kappa$或$\lambda$无限, 那么
    $\kappa\otimes\lambda=\max\{\kappa,\lambda\}$.
\end{enumerate}

\noindent{\kaishu 证明.}

\begin{enumerate}
    \item 如果$\kappa,\lambda$都有限, 那么$\kappa\lambda$是自然数,
    因此$\kappa\otimes\lambda=|\kappa\lambda|=\kappa\lambda$.
    \item 如果$\kappa$或$\lambda$无限且$\kappa\geqslant\lambda$,
    那么$\kappa$无限. 一方面
    \[
        \kappa\lambda\preccurlyeq\kappa\kappa\approx
        \kappa\times\kappa\approx\kappa,
    \]
    另一方面$\kappa\lambda\geqslant\kappa$, 所以$\kappa\lambda\approx\kappa$,
    即$\kappa\otimes\lambda=\kappa$.
    
    \item 如果$\kappa$或$\lambda$无限且$\kappa\leqslant\lambda$,
    则同理得$\lambda\kappa\approx\lambda$, 所以$\kappa\otimes\lambda=\lambda$.
    \hfill $\qedsymbol$
\end{enumerate}

\vspace{10pt}

\hypertarget{sec:定理1.3}{}
\noindent\textbf{定理1.3}

\begin{enumerate}
    \item 结合律: $\forall\kappa,\lambda,\mu\in\mathbf{Cn}:\,
    (\kappa\otimes\lambda)\otimes\mu=\kappa\otimes(\lambda\otimes\mu)$,
    \item 零元: $\forall\kappa\in\mathbf{Cn}:\,\kappa\otimes 0=0$,
    \item 单位元: $\forall\kappa\in\mathbf{Cn}:\,\kappa\otimes 1=\kappa$,
    \item 交换律: $\forall\kappa,\lambda\in\mathbf{Cn}:\,
    \kappa\otimes\lambda=\lambda\otimes\kappa$,
    \item 分配律: $\forall\kappa,\lambda,\mu\in\mathbf{Cn}:\,
    \kappa\otimes(\lambda\oplus\mu)=(\kappa\otimes\lambda)\oplus(\kappa\otimes\mu)$.
\end{enumerate}

\end{document}